\documentclass[11pt]{article}
\usepackage[a4paper,margin=19mm]{geometry}
\usepackage{fontspec}
\setmainfont{lmroman10-regular.otf}
\usepackage{amsmath,amssymb,amsthm,mathtools,amscd,xcolor,graphicx,hyperref}
\hypersetup{hidelinks}
\numberwithin{equation}{section}
\newtheorem{theorem}{Theorem}[section]
\newtheorem{lemma}[theorem]{Lemma}
\newtheorem{proposition}[theorem]{Proposition}
\newtheorem{corollary}[theorem]{Corollary}
\newtheorem{claim}[theorem]{Claim}
\newtheorem{fact}[theorem]{Fact}
\newtheorem{conjecture}[theorem]{Conjecture}
\theoremstyle{definition}
\newtheorem{definition}[theorem]{Definition}
\newtheorem{example}[theorem]{Example}
\newtheorem{remark}[theorem]{Remark}
\newtheorem{question}[theorem]{Question}
\theoremstyle{plain}
\newcommand{\C}{\mathbb C}
\newcommand{\R}{\mathbb R}
\newcommand{\Q}{\mathbb Q}
\newcommand{\Z}{\mathbb Z}
\newcommand{\N}{\mathbb N}
\newcommand{\eee}{\mathrm e}
\newcommand{\cclass}[1]{\textsf{#1}}
\newcommand{\cproblem}[1]{\textsc{#1}}
\newcommand{\ie}{i.\,e.}
\newcommand{\Ie}{I.\,e.}
\newcommand{\eg}{e.\,g.}
\newcommand{\Eg}{E.\,g.}
\newcommand{\phd}{Ph.\,D.}
\newcommand{\msc}{M.\,S.}
\newcommand{\bsc}{B.\,S.}
\newcommand{\expref}[2]{\hyperref[#2]{#1~\ref{#2}}}
\newcommand{\expeqref}[2]{\hyperref[#2]{#1~\eqref{#2}}}
\newcommand{\exprefsilent}[2]{\hyperref[#2]{#1}}
\providecommand{\etal}{et~al.}
\usepackage{eprint}
\usepackage[nameinlink,capitalize]{cleveref}
\renewcommand{\eqref}[1]{\hyperref[#1]{(\ref*{#1})}}

\usepackage{thmtools, thm-restate}
\usepackage{bbm}

\theoremstyle{plain}
\newtheorem{mainthm}[theorem]{Main Theorem}
\newtheorem{mcorollary}[theorem]{Main Corollary}
%

\DeclareMathOperator{\poly}{poly} %
\DeclareMathOperator{\dist}{dist} %
\DeclareMathOperator{\mindist}{mindist}  %
%
%
%
%

\newcommand{\K}{{\mathbb{K}}}
\newcommand{\E}{{\mathbb{E}}}
\newcommand{\F}{{\mathbb{F}}}
%
%
\newcommand{\calT}{{\mathcal{T}}}
\newcommand{\calS}{{\mathcal{S}}}
\newcommand{\calK}{{\mathcal{K}}}
\newcommand{\calL}{{\mathcal{L}}}
\newcommand{\calI}{{\mathcal{I}}}


\newcommand{\eps}{\epsilon}
%
%
%
\newcommand{\plurality}{{\rm{plurality}}}
%

\newcommand{\func}[1]{\{\F_q^{#1} \to \F_q\}}
\newcommand{\Func}[1]{\{\F_q^{#1} \to \F_q\}}
\newcommand{\pl}{{\rm{Plur}}}
\newcommand\sett[2]{\left\{ #1 \left| \; \vphantom{#1 #2} \right. #2  \right\}}



%
%
%
\title{Affine-invariant lifted codes: local testing implies robust testing through agreement}
\author{Irit Dinur, Prahladh Harsha, Tali Kaufman, Noga Ron-Zewi}
\date{doi:10.4086/toc.2022.v018a012; CC BY 3.0}
\begin{document}
\maketitle
\noindent\textit{Editorial note.} Complete original mathematical paper, all supporting author proofs and original published bibliography retained. Native frontmatter/biographies remain in raw source/publication PDF. Existing article interface adapts XeLaTeX presentation only.\par\medskip
\noindent\textit{Editorial note.} One primary construction family is counted by this task. Related consequences, special cases and restatements remain complete supporting context. Original theorem/environment names and all exact assumptions retained.\par\medskip
\noindent\textit{Editorial note.} Exact original CC BY3.0 retained. Authors own mathematics; no proof review, mathematical repair, added assumptions, content hashes or cross-version comparison.\par\medskip
\noindent\textit{Allowed standard prerequisites (editorial).}\begin{itemize}
\item Original dual-polynomial or affine-code definitions and cited standard tools. All original degree/error/family/depth or field/dimension/minimum-distance/soundness assumptions retained literally.. Complete source-local correction/amplification or sampling/agreement construction supplied, without collector reproof or certification.
\end{itemize}\medskip
\section{Introduction}

Our main result shows a transformation from \emph{local testing} to \emph{robust testing} for the class of \emph{affine-invariant lifted codes}. We start by describing the notions of local testing, robust testing, and lifted codes.

\subsection{Local testing and robust testing}

A code is a subset $C \subseteq \Sigma^n$. The elements of
%
%
%
$\Sigma^n$ are called \emph{words}, the elements of %
$C$ are called \emph{codewords}, $\Sigma$ is the \emph{alphabet} of the code, and $n$ is the \emph{block length}. The \emph{rate} of the code is the ratio $(\log_{|\Sigma|}|C|) / n$. The code is \emph{linear} if $\Sigma = \F_q$ where $\F_q$ is the finite field of $q$ elements, and $C$ is an $\F_q$-linear subspace of $\F_q^n$.
It will be convenient to think of codewords in $C$ as functions $f: U \to \Sigma$ where $U$ is a domain of size $n$.
For a pair of functions $f,g: U \to \Sigma$, we let $\dist(f,g)$ denote
%
their normalized Hamming distance, \ie,  %
the fraction of inputs  $x \in U$ for which $f(x) \ne g(x)$.
%
%
By the \emph{minimum distance}  %
$\mindist(C)$ of the code  %
we mean its minimum normalized distance, \ie,  %
%
the minimum of $\dist(f,g)$ over all
pairs of %
distinct codewords $f,g \in C$.
For a function $f: U \to \Sigma$, we let  $\dist(f,C)$ denote
the minimum of $\dist(f,g)$ over all codewords $g \in C$.

%
%
%
Let $Q\in\Z^+$ and $\alpha > 0$.   %
A \emph{local tester} for the code $C$
with \emph{query complexity $Q$} and \emph{soundness $\alpha$}  %
is a probabilistic oracle algorithm that
on oracle access to a function $f:U \to \Sigma$
makes at most $Q$ queries to $f$, and accepts $f \in C$
with probability one
\emph{(completeness)}, %
%
%
while rejecting %
%
all words $f\in\Sigma^n$  %
with
probability at least $\alpha \cdot \dist(f,C)$.
%
%
In this %
article,  %
we shall restrict our attention to local testers that pick a random subset 
 $K \subseteq U$ of cardinality $Q$ according to some distribution, and 
accept%
%
\footnote{Local testers may generally apply a more complex predicate on $f|_K$. However, natural local testers are typically of the restricted form we consider, and moreover it can be shown that a local tester for a linear code
must be of this form \cite{BenSassonHR2005}.}
if and only if $f|_K \in C|_K$.   %
%
%
%
%
%
Then the completeness requirement is automatically satisfied
($f|_K \in C|_K$ with probability one whenever $f \in C$),
and we only need to consider the soundeness condition,   %
\begin{equation}\label{eq:test}
\Pr_{K}[f|_K \notin C|_K ] \geq \alpha \cdot \dist(f,C)
\end{equation}
%
for all words $f$.  %

In this %
article  %
we will be interested in the stronger notion of robust soundness.\footnote{This strengthened notion of soundness has been referred to as \emph{robust soundness} \cite{BenSassonGHSV2006,BenSassonS2006} or just \emph{robustness} \cite{GuoHS2015} in earlier
%
work.  %
To make precise the contrast with the standard notion of soundness \eqref{eq:test}, we will refer to %
it  %
as ``robust soundness''  in this
%
paper.}   %
We say that a local tester
%
%
%
is \emph{robustly sound} (or just \emph{robust}) if
%
for all words,   %
the average distance of its local views from accepting views is
at least  %
%
proportional to the distance of the given function from the code.
That is, 
%
%
%
%
%
instead of \eqref{eq:test} we now require that
\begin{equation}\label{eq:robust}
\E_K\big[\dist(f|_K,C|_K)\big] \geq \alpha \cdot \dist(f,C)
\end{equation}
%
for all words $f$.  %
%
%
We   %
refer to $\alpha$ as the \emph{robust soundness} 
of the tester.

The notion of robust soundness was introduced by Ben-Sasson and Sudan \cite{BenSassonS2006} based on analogous notions for probabilistically checkable proofs \cite{BenSassonGHSV2006,DinurR2006}.
Robust %
soundness   %
is a natural property of local testers that relates the global distance of a function from the code to its local distance from accepting views on local views. Moreover, robust soundness is also an important ingredient in constructions of locally testable codes and probabilistically checkable proofs as it allows for composition of local tests. Specifically, it follows by definition that if a code $C$ is robustly testable with query complexity $Q$ and soundness $\alpha$, and additionally each local restriction $C|_K$ is locally testable with query complexity $Q'$ and soundness $\alpha'$, then the code $C$ is locally testable with query complexity $Q'$ and soundness $\alpha \cdot \alpha'$. This property is useful when local restrictions can be tested efficiently, %
which can happen if the code has many symmetries (as is the case with
the class of lifted codes considered in this %
paper),  %
or can be achieved, in the case of probabilistically checkable
%
proofs,  %
by attaching a short proof of proximity.

One can easily observe that  \eqref{eq:robust} implies \eqref{eq:test} since $f|_K \notin C|_K$ whenever $\dist(f|_K,C|_K)>0$, so robust soundness is a stronger requirement than (standard) soundness in local testing. 
For the other direction, note that a local tester with soundness $\alpha$ has robust soundness at least $\alpha/Q$ since
$\dist(f|_K,C|_K) \geq 1/Q$ whenever $f|_K \notin C|_K$. A natural question is whether this loss in robust soundness is necessary, and whether robust soundness is
a  %
strictly stronger notion than (standard) soundness for local testing.

In this %
paper %
we shall show that this loss is unnecessary for the class of lifted codes, discussed below.

\subsection{Lifted codes}

%
%
For sets $B,\Sigma$ let $\{B \to \Sigma\}$ denote the
set of all functions that map $B$ to $\Sigma$.
We can view this set as the set of strings of
length $n=|B|$ over $\Sigma$.
If $\Sigma=\F_q$ then this set is a vector space over
$\F_q$ of dimension $|B|$.   %

Lifted codes are specified by a \emph{base code} 
$C \subseteq \{\F_q^\ell \to \F_q\}$ and a dimension $m \geq \ell$.
We further assume that the base code $C$ is linear and
%
%
affine invariant.   %
That is, for any codeword $f \in C$, 
and for any affine transformation $A: \F_q^\ell \to \F_q^\ell$,
it holds that $f \circ A \in C$.
Given these we define the \emph{lifted code} $C^{\ell \nearrow m}$ to be the code consisting of all functions $f: \F_q^m \to \F_q$ that satisfy that $f|_L \in C$ for any $\ell$-dimensional affine subspace $L$
of the domain $\F_q^m$.  %
%

Lifted codes were first introduced by Ben-Sasson, Maatouk, Shpilka and Sudan \cite{BenSassonMSS2011}, and their local testability properties were further explored in subsequent work  \cite{GuoKS2013,HaramatyRS2015,GuoHS2015}.
They are a natural generalization of the well-studied family of Reed--Muller codes, and moreover they also give rise to new families of locally testable codes that outperform Reed--Muller codes in a certain range of parameters \cite{GuoKS2013}. Specifically, lifted codes lead to one of the two known constructions (the other one being tensor codes \cite{BenSassonS2006,BenSassonV2015,Viderman2015,KoppartyMRS2017}) of high-rate locally testable codes (\ie, locally testable codes with rate approaching one and
%
%
locality $n^{\beta}$ for any constant $\beta>0$). 
%
Generally, lifted codes
form a natural subclass of affine-invariant codes %
%
having  %
the `single-orbit characterization' %
property that is known to imply local testability, as well as local decodability \cite{KaufmanS2008}.

There is a natural local test associated with lifted codes: on oracle access to a function $f: \F_q^m \to \F_q$, pick a uniform random $\ell$-dimensional affine subspace %
$L\subseteq \F_q^m$  %
and accept if and only if $f|_L \in C$. It follows immediately by definition that this test accepts any valid codeword $f \in C^{\ell \nearrow m}$ with probability one, but more work is required to show that this test is sound. Specifically, since the test forms a
%
single-orbit   %
characterization, it follows from  the work of Kaufman and Sudan \cite{KaufmanS2008} that it has soundness roughly $q^{- 2\ell}$. The dependence of the soundness on the dimension $\ell$ was later eliminated by Haramaty, Ron-Zewi and Sudan \cite{HaramatyRS2015} who showed soundness that is only a function of $q$ (though an extremely quickly decaying one).

As for robust soundness, the above local testing results, together with the straightforward transformation from local testing to robust testing, immediately give robust soundness that is dependent on the dimension $\ell$. This
dependence %
was eliminated recently by Guo, Haramaty and Sudan \cite{GuoHS2015} who showed robust soundness of the form $\poly(\delta)$ (about $\delta^{74}$,  where $\delta$ is the %
minimum     %
distance of the code) for the local test that queries subspaces of slightly larger dimension of $2\ell$. Interestingly,  Guo, Haramaty and Sudan \cite{GuoHS2015} did not rely on the aforementioned local testing results, but rather relied on viewing lifted codes as the intersection of
`modified tensor codes.'  %
They then proceeded by showing that these modified tensor codes are robustly testable (using %
Viderman's method   %
\cite{Viderman2015} showing robust soundness of tensor codes), and that this implies local testability of the lifted
%
code.  (See  %
\cref{subsec:proof-prev} for more details about the 
%
method of  Guo, %
Haramaty  %
and Sudan \cite{GuoHS2015}.) %

\subsection{Our results}

Our main result gives a transformation from local testing to robust testing,  that does not suffer the factor of $Q$ (the query complexity) loss in robust soundness, for the class of lifted codes. The transformation uses local testability in a `black-box' manner, and shows that if a code in this family is locally testable (using the natural subspace tester) then it is also robustly testable with roughly the same number of queries and robust soundness.

For $k \geq \ell$, let the
%
%
%
%
%
\emph{$k$-dimensional test}
denote the local tester that on oracle access to a function 
$f: \F_q^m \to \F_q$ queries a uniform random $k$-dimensional 
affine subspace
%
$K \subseteq \F_q^m$  %
and accepts if and only if $f|_K \in C^{\ell \nearrow k}$.


\begin{restatable}{mainthm}{maint}\label{thm:main}
%
Let $C \subseteq \Func{\ell}$ be an affine-invariant linear code, and
$m \geq k \geq \ell$. Suppose that $C^{\ell \nearrow m}$ is locally 
testable using the
$k$-dimensional test with query complexity $q^{k}$ and soundness $\alpha$, 
and let  $\delta:= \min_{ k \leq r \leq m}\mindist(C^{\ell \nearrow r})$.
Then $C^{\ell \nearrow m}$ is robustly testable using the $(2k+2+\log_q(4/\delta))$-dimensional test
 with query complexity $O(q^{2k+2}/\delta)$ and
  robust soundness $\Omega(\alpha \cdot \delta^3)$.
%
\end{restatable}

  
Note that if the %
minimum  %
distance $\delta$ is constant, we only incur a constant multiplicative loss in robust soundness and testing dimension.\footnote{One natural example for codes with constant $\delta$ are the lifted Reed-Solomon codes of Guo, Kopparty and Sudan \cite{GuoKS2013}. Furthermore, it was shown by them \cite{GuoKS2013} that for any $m \geq \ell$, $\mindist(C^{\ell \nearrow m})\geq\mindist(C)-q^{-\ell}$ (see \cref{lifted-dist} below).} We conjecture that the testing dimension for proving robust soundness may be as small as $k+1$, and leave it as an interesting question for future research.

To apply the above theorem one can instantiate it with the local testing result of Kaufman and Sudan \cite{KaufmanS2008}  that says that lifted codes are locally testable using the $\ell$-dimensional test with soundness $\approx q^{-2\ell}$ (see \cref{lifted-testing} below).
However, to obtain constant robust soundness we need that the soundness 
of the initial local tester
%
be constant (independent of $q$ and $\ell$), and for this we observe (in \cref{lifted-testing-ampl}) that the soundness of Kaufman and Sudan \cite{KaufmanS2008} can be easily amplified to $\Omega(1)$ at the cost of increasing the testing dimension%
%
%
\footnote{Such an amplification with
a  %
similar blow-up in query complexity can be easily obtained by repeating the test and accepting if and only if all invocations accept; we,
 however,  %
 need that the tester %
 be a subspace tester which can be obtained using sampling properties of affine subspaces.}   
to $\approx 3\ell$. %
Using this observation we obtain the following.

\begin{restatable}{mcorollary}{maincor}\label{cor:main}
%
There exists an absolute constant $c>1$ so that the following holds.
Let $C \subseteq \Func{\ell}$ be an affine-invariant linear  code of
%
minimum   %
distance $\delta$, and $m \geq \ell$. Then $C^{\ell \nearrow
  m}$ is robustly testable using the
$(6\ell+4+\log_q(c/\delta))$-dimensional test  with robust soundness
$\Omega( \delta^3)$.
%
\end{restatable}


Compared to the above corollary, the %
Guo, Haramaty and Sudan \cite{GuoHS2015} %
use  %
a lower dimension of $2\ell$, but also obtain a 
lower robust soundness of $\Omega(\delta^{74})$.

As described next, our proof is combinatorial, relying mainly on sampling properties of the collection of local views. In particular, it uses very little about the algebraic structure of lifted codes or the base code.
We thus hope that such techniques %
may  %
prove useful in the future for showing robust soundness for 
other families of locally testable codes with similar sampling properties.

%
\section{Overview of the proof}  %
\label{sec:proof-overview}

Our proof is based on a new connection between the notions of local testing, robust testing, and \emph{agreement testing}. Specifically, we show that for the class of lifted codes, %
 agreement testing implies robust testing, and local testing implies agreement testing. The combination of these two implications gives our \cref{thm:main}. Next we elaborate on the notion of agreement testing, followed by an overview of each of the implications.

\subsection{Agreement testing}

An \emph{agreement test} attempts to find out whether
%
%
assignments to local views 
can be stitched together to a single global codeword. Let $C \subseteq \{U \to \Sigma\}$  be a code, and let $\calS$ be a collection of subsets of $U$.
An \emph{agreement tester} for %
$(C, \calS)$   %
%
is a probabilistic oracle algorithm that receives oracle access to a collection of partial assignments $\{f_S:S\to \Sigma \,\mid\, S\in \calS\}$,
%
where $f_S \in C|_S$ for any $S \in \calS$. The tester queries a few of 
%
%
the $f_S$ %
%
and is required to accept with probability one any collection $(f_S)_S$ that is consistent with some global codeword $g \in C$ (that is, $g|_S = f_S$ for any $S \in \calS$), while rejecting any inconsistent collection $(f_S)_S$
with probability
at least  %
proportional to the minimal fraction
%
of the $f_S$  %
that  must be changed in order to be consistent with some global codeword. In this %
paper,   %
we focus on an agreement tester that picks a pair of sets $S,S' \in \calS$ according to some distribution and accepts if and only if $f_S$ and $f_{S'}$ agree on their intersection $S\cap S'$.

Agreement testing has first appeared in PCP constructions  \cite{BabaiFL1991,BabaiFLS1991,FeigeGLSS1996,AroraS1998,AroraLMSS1998}  as so-called ``low-degree tests,'' and is a key component in the proof of almost all PCP theorems.
A prime example is the line vs.\ line low-degree test  \cite{GemmellLRSW1991,RubinfeldS1996} in the proof of the PCP theorem. In the PCP construction, a function on a large vector space is replaced by an ensemble of (supposed) restrictions to all possible affine lines. These restrictions are supplied by a prover and are not a~priori guaranteed to agree with any single global function. The  ``low-degree test'' is run by the verifier to check that restrictions on intersecting lines agree with each other, \ie, they give the same value to the point of intersection. The main point of the argument is to show that the passing of the test implies agreement with a single global function.
In these early low-degree tests (including the linearity
%
tester  %
of Blum, Luby and Rubinfeld \cite{BlumLR1993}) an agreement test component can be discerned but quite implicitly. Indeed, it was only separated in the
%
work of Raz, Safra, Arora, and Sudan %
%
\cite{RazS1997,AroraS2003} that looked at the so-called list-decoding regime,\footnote{In the list decoding regime one would like to reject a function that is $(1-\epsilon)$-far from the code with very
%
high probability  %
of $1-O(\epsilon)$.} with the goal of proving a large gap for the PCP.

Goldreich and Safra \cite{GoldreichS2000} tried to separate the algebraic aspect of the low-degree test from the combinatorial, and formulated a more general ``consistency test'' which is also referred to as an agreement test. They also proved a certain local to global result which was too weak to be useful for PCPs. In hindsight it is clear that since their family of subsets consisted of
%
axis-parallel  %
lines, the expansion was not strong enough for a good agreement test.
Only recently \cite{DinurK2017, DiksteinD2019}
%
has %
the role of expansion underlying the family of subsets %
begun to be uncovered. 

Work on agreement testing then continued the combinatorial direction of Goldreich and Safra~\cite{GoldreichS2000} mainly in the list-decoding regime
for direct product testing
\cite{DinurR2006,DinurG2008,ImpagliazzoKW2012,DinurS2014-dp,DinurN2017}. The techniques developed in this line of work turn out to be useful also for agreement testing in the unique-decoding regime (which is the more standard testing regime), and in particular for our work here.

Our proof of agreement testing based on local testing (see \cref{sec:testing-to-agreement}) is significantly simpler than general agreement testing theorems. The reason is that we avoid a major technical difficulty that stems from the varying size of disagreement between pairs of local views. In the general setting this can be anywhere from disagreement on only one point in the intersection to disagreement on the entire intersection, leading to a very subtle argument. In contrast, in our setting the local view restricted to the intersection is itself an error correcting code, so whenever  local views disagree, they must disagree on a %
constant  %
fraction of the intersection. This makes the proof of the agreement theorem %
quite  %
direct, and much easier than the proofs in
the  %
aforementioned %
papers.  %

%
%
An agreement  %
testing theorem that holds in quite broad generality 
was proven by Dikstein and Dinur \cite{DiksteinD2019} 
(published after the %
conference version of the present paper~\cite{conf-version}).
%
%
This result (see Theorems 6.2 and 6.3 in %
\cite{DiksteinD2019}) is   %
%
similar to what we prove here but the parameters of Dikstein and Dinur are slightly weaker.
%
%

\subsection{Agreement testing implies robust testing}

We begin with an overview of the simpler implication from agreement
testing to robust testing. In what follows, we say that a
collection %
$\calS$ of subsets %
of the ground set $U$ \emph{samples well}
%
%
%
%
%
%
\emph{inside $U$,}   %
if, informally,
for all $0\le \gamma\le 1$ and %
%
%
for any subset $A \subseteq U$ of density $\gamma$ (\ie, $|A|/|U| = \gamma$),
it holds that for most subsets $S \in \calS$, the density of $A$ inside $S$
is approximately $\gamma$
%
(\ie, $|A\cap S| /|S| \approx \gamma$). %
%

Suppose that we have an agreement tester for  %
$(C, \calS)$   %
as described above. We would like to show that the local tester that 
queries a random $S \in \calS$ is robustly sound.
Let $\calT$ be the collection of subsets of $U$ formed by
the %
%
pairwise intersections of
distinct  %
%
sets in $\calS$.
%
%
%
The main properties
%
of $\calS, \calT$ we need    %
are sampling properties, 
specifically, that $\calS$ samples well
%
inside $U$   %
and that for any $S \in \calS$, %
the sets  %
in $\calT$ contained in $S$ sample well 
%
inside $S$.   %
 The main property %
of the code $C$
we need  %
is that for any set $T \in \calT$, the code $C|_T:=\{c|_T \mid c\in C\}$
has large distance.  %
In the case of lifted codes these properties can be guaranteed by letting 
%
$\calS$ and $\calT$  %
be families of affine subspaces of fixed dimensions
%
$s$ and $t$,  %
respectively, where%
%
%
%
%
\footnote{In this paper, we use the informal notation $s \gg t$
for positive numbers $s,t$ to indicate $s$ is \emph{much greater}
than $t$.}   %
$s \gg t$.

To see that the proposed local tester is indeed robustly sound, suppose
we have a function $f: U \to \Sigma$ that is close to  $C|_S$ on a
%
typical $S\in\calS$.  %
Our goal is to show that $f$ is close to a codeword $g \in C$. We first create an instance $(f_S)_S$ for the agreement tester by letting $f_S \in C|_S$ be the closest valid assignment to $f|_S$. Next observe that since $f|_S$ is typically close to $f_S$, and by assumption 
%
the sets $T$ sample well inside the sets $S$,  %
for a typical $T$ and $S,S'$ containing $T$, it holds that both $f_S|_T$ 
and  $f_{S'}|_T$ are close to $f|_T$, and by the
%
property that $C|_T$ has large distance  %
%
%
for any $T \in \calT$,
this implies in turn that typically $f_S|_T = f_{S'}|_T$.
Consequently, agreement testability implies the existence of a codeword $g \in C$ that agrees with most $f_S$, and so $g|_S= f_S$ and $f_S$ is close to  $f|_S$ for most $S$. But since $\calS$ samples well inside $U$ we conclude that $f$ must be close to $g$ as required.

\subsection{Local testing implies agreement testing}

We now turn to the local testing to agreement testing implication which is a bit more involved. Suppose that we have a local testing algorithm for $C$ that queries a random set $K\in \calK$ and accepts if and only if $f|_K \in C|_K$. We would like to obtain an agreement tester for $C$ with respect to some collection of subsets $\calS$.
As before, let $\calT$ be the collection of subsets of $U$ formed by 
the  %
%
pairwise intersections of sets in $\calS$.
Once more the main properties 
%
of $\calS, \calT, \calK$
we require  %
are sampling properties. Specifically, we need  that $\calS$ samples well inside $U$, and that for any $T \in \calT$, all sets in $\calK$ contained in $T$ sample well inside $T$.
We also require distance properties %
of $C$, specifically that 
$C, C|_T, C|_S$ all have large distance %
for any $S \in \calS$ and $T \in \calT$.
 Once more, in the case of lifted codes these properties can be guaranteed by letting $\calS, \calT, \calK$ be families of affine subspaces of fixed dimensions $s,t,k$, respectively, where $s \gg t \gg k$.

To show agreement testability, let  $(f_S)_S$ be a collection of valid assignments to sets in $\calS$ (so  $f_S \in C|_S$ for any $S$), and suppose that $f_S$ agrees with $f_{S'}$ on $S\cap S'$ for most pairs $S,S'$. Our goal will be to find a global codeword $g \in C$ that agrees with most $f_S$.
We find the function $g$ in the following three stages.

\paragraph{Initial stage.} In the first stage we define for any $K \in \calK$ a `most popular function' %
$\pl_K$ by choosing the most common value among $f_S|_K$ going over all $S \in \calS$ containing $K$. We then show, using the assumption
that %
the $f_S$  %
typically agree on their
pairwise  %
intersections, that this most popular function is obtained with 
overwhelming probability for a typical $K$.
%
%

\paragraph{Local structure stage.} In the second stage we define for each $K \in \calK$ a function $g_K: U \to \Sigma$ by letting $g_K(x)$ be the most common `vote' among all $f_S$
%
where $S\supseteq K$ and $x\in S$  %
and %
$f_S$ agrees  %
with $\pl_K$  on $K$ (this function is well defined because of the initial stage). We then show that for a typical $K$, $g_K$ is close to some function $h_K \in C$, and moreover $h_K|_S = f_S$ for most $S$ containing $K$.

To see why the above holds, first note that by the assumptions 
that $C|_T$ has large distance %
for any $T \in \calT$, 
 %
 and the sets $K$  %
sample well inside %
each $T$  %
%
if a pair of %
partial functions $f_S$ and $f_{S'}$  %
agree on $K$ then they must typically also agree on their whole intersection. Therefore $g_K(x)$ is also typically defined with overwhelming probability. Consequently, for a typical $K$, and most $K'$,\ $g_K|_{K'}$ agrees with some $f_S$. Recalling that
%
the $f_S$ are  %
valid assignments, local testability then implies the existence of $h_K \in C$ that is close to $g_K$. The fact that $h_K|_S = f_S$ for most $S$ containing
$K$ follows by the assumption that $\calS$ samples well
inside  %
$U$, and the distance property on $\calS$. 

\paragraph{Global structure stage.} In the final stage we show that there 
exists %
${\hat K}\in \calK$  %
such that $h_{\hat K}$ agrees with $f_S$ for most $S$ (not necessarily 
containing $\hat K$). We can then set our
`global function' %
$g$ to be equal to $h_{\hat K}$.
To this end, we first observe that it suffices to show that most functions
$h_K$ are in fact identical.
This now follows since for typical $S \supseteq K\cup K'$ it holds that 
$h_K|_S = f_S = h_{K'}|_S$, and consequently since $\calS$
samples %
well inside $U$,  %
it must typically hold that $h_K=h_{K'}$.
%
%

\subsection{The proof method of Guo, Haramaty and Sudan}\label{subsec:proof-prev}

%
Guo, Haramaty and Sudan's method \cite{GuoHS2015} %
%
to show  %
robust soundness of lifted codes is very different from ours. In particular, it relies heavily on the algebraic structure of lifted codes. More specifically, the proof is based on viewing the lifted codes as the intersection of
`modified tensor codes.'  %
%
%
The \emph{tensor power}   %
$C^{\otimes m}$ of a code $C \subseteq \{ \F_q \to \F_q\}$ can be 
thought of as the
`axis-parallel lifting'  %
of $C$. That is, it is the code that consists of all functions
 $f: \F_q^m \to \F_q$ whose restrictions to
any axis-parallel line belong to $C$. The
`modified tensor code'   %
is a code of the form $C^{\otimes m}_{b}$ where $b$ is a direction in $\F_q^m$,
and $C^{\otimes m}_{b}$ consists of all functions $f \in C^{\otimes m}$ whose 
restrictions to lines in direction $b$  also belong to $C$.

%
%
The authors first use Viderman's method \cite{Viderman2015}, %
showing robust testing of tensor codes, to show that the modified tensor codes are also robustly testable. They then use the fact that the lifted code is the intersection of all codes of the form $C^{\otimes m}_{b}$ for all directions $b$ (this is true when the dimension of the base code for lifting is $\ell=1$; when $\ell >1$ the proof becomes more complicated) to deduce robust testability for the lifted code. However, since
the  %
intersection of robustly testable codes is not necessarily robustly testable, %
non-trivial work is required to show robust testability, which in particular exploits the degree structure of affine-invariant lifted codes.

The above program can be carried out only when the dimension $m$ of the lifted code is a small constant multiple of $\ell$, and the authors use the
%
%
%
bootstrapping technique of
\cite{RubinfeldS1996,AroraLMSS1998,AroraS2003,Arora1994} 
to extend the result to
%
arbitrarily   %
large $m$.

In contrast, we work directly with lifted codes of large dimension which allows us to exploit the sampling  properties of large affine subspaces in $\F_q^m$. To the best of our knowledge, even for the special case of low-degree polynomials, this gives the first
analysis of robust soundness that  is not
based on the %
two-step   %
approach of first analyzing the
constant-dimensional %
case and only then moving to the general
%
case.

In contrast to the proof of Guo, Haramaty and Sudan \cite{GuoHS2015} 
who reprove  local testability %
along %
the way, our proof uses local testability in a black-box manner. Thus, it exhibits a separation between the algebraic properties that are used for showing local testability, and the combinatorial properties that  are needed in order to turn local testability into robust testability.

%
%
%
%

%
\paragraph{Organization of the paper.}  %
The rest of the paper is organized as follows. In \cref{sec:prelim} we set some notation, provide some definitions, and present the sampling properties of subspaces that we use. The transformation from agreement testing to robust testing is given in \cref{sec:agreement-to-robust}, while the transformation from local testing to agreement testing
%
appears  %
in \cref{sec:testing-to-agreement}.
We %
%
%
conclude in %
\cref{sec:testing-to-robust} with the full transformation
from local testing to robust testing that proves 
our \cref{thm:main} and \cref{cor:main}.


\section{Preliminaries}\label{sec:prelim}

%
%
%
%
For a prime power $q$, let $\F_q$ denote the finite field of $q$ elements.
Let $\Func{m}$
denote the set of functions mapping $\F_q^m$ to  $\F_q$. In what follows
we focus on codes which are
%
sets  %
of functions, %
$C \subseteq \Func{m}$. For a pair of functions $f,g: \F_q^m \to \F_q$,
we use $\dist(f,g)$ to denote the
fraction of inputs $x \in \F_q^m$ for which $f(x) \ne g(x)$.
The  %
minimum distance   %
$\mindist(C)$ of the code $C$ is $\min_{f \neq g\in C}
\{\dist(f,g)\}$. For a function $f: \F_q^m \to \F_q$,
we use $\dist(f,C)$ to denote $\min_{g \in C}
\{\dist(f,g)\}$.  %

The code $C$ is said to be \emph{linear}
if it is an $\F_q$-linear subspace, 
%
%
%
\ie, $C\neq\emptyset$ and %
for every $\alpha \in \F_q$ and $f,g \in C$, we
have $\alpha f + g \in C$.   %
A function $A:\F_q^m \to \F_q^m$
is said to be an \emph{affine transformation} if there exist a matrix $M \in \F_q^{m \times m}$
and  a vector $b \in \F_q^m$ such that $A(x) = Mx + b$.
The code $C$ is said to be \emph{affine invariant} %
if for every affine transformation $A$ and every $f \in C$
we have $f \circ A \in C$ (where $(f \circ A)(x) = f(A(x))$).

\subsection{Lifted codes}

A subset $L \subseteq \F_q^m$
is said to be an \emph{$\ell$-dimensional affine subspace} if there
exist
%
%
%
%
$u_0 \in \F_q^m$ and linearly independent $u_1,\ldots,u_\ell \in \F_q^m$ such that
%
$L = \{u_0 + \sum_{i=1}^\ell u_i x_i \mid x_1,\ldots,x_\ell \in \F_q\}$.   %
%
%
%
We fix an arbitrary %
invertible  %
affine map $\gamma_L:\F_q^\ell\to L$ (which we can view as a parameterization of $L$). For a function $f:\F_q^m \to \F_q$, the restriction $f|_L$ is viewed as a function in $\{\F_q^\ell \to \F_q \}$ through $f\circ\gamma_L: \F_q^\ell\to\F_q$. In particular, when we ask if $f|_L \stackrel{?}\in C\subseteq \{ \F_q^\ell \to \F_q \}$ what we are really asking is whether $f\circ \gamma_L \in C$. Note that if $C$ is affine-invariant, whether $f|_L \in C$ does not depend on the choice of the parametrization $\gamma_L$.
\begin{definition}[Lifted codes]
Let $C \subseteq \Func{\ell}$ be an affine-invariant  linear code, and $m \geq \ell$. The \emph{$m$-dimensional
lift} $C^{\ell \nearrow m}$ of $C$ is given by
\[C^{\ell \nearrow m}:=\big\{f :\F_q^m \to \F_q \mid f|_L \in C \; \text{for every $\ell$-dimensional
affine subspace $L \subseteq \F_q^m$}\big\}\,.\]
\end{definition}


\begin{proposition}[Distance of lifted codes, {\cite[Theorem 5.1, Part (2)]{GuoKS2013}}] \label[proposition]{lifted-dist}
Let $C \subseteq\{\F_{q}^{\ell}\to\F_{q}\}$ be an affine-invariant  linear code, and $m \geq \ell$. Then $\mindist(C^{\ell \nearrow m})\geq\mindist(C)-q^{-\ell}$.
\end{proposition}


\subsection{Local testing, robust testing, and agreement testing}

We now formally define the notions of local testing, robust testing, and agreement testing, specialized to the class of lifted codes and subspace testers. In the case of local testing and robust testing this simply means that the tester samples a uniform random $k$-dimensional affine subspace and its
accepting views are codewords in $C^{\ell \nearrow k}$.

\begin{definition}[Local testing of lifted codes]
Let $C \subseteq \Func{\ell}$ be an affine-invariant linear  code, and $m\geq k \geq \ell$. The $m$-dimensional lift $C^{\ell \nearrow m}$ is \emph{$(k,\alpha)$-testable} if for every $f: \F_q^m \to \F_q$ it holds that
\[\Pr_K\big[f|_K \notin C^{\ell \nearrow k}\big] \geq \alpha \cdot \dist(f, C^{\ell \nearrow m})\,,\]
where the probability is over a uniform random $k$-dimensional affine subspace $K \subseteq \F_q^m$.
\end{definition}

%
%
Recall the definition of the \emph{$\ell$-dimensional test}
in the paragraph before \cref{thm:main}.   %

\begin{theorem}[{\cite[Theorem 2.9]{KaufmanS2008}}]\label{lifted-testing}
  Let $C \subseteq\{\F_{q}^{\ell}\to\F_{q}\}$ be an affine-invariant linear  code, and $m \geq \ell$. Then the $\ell$-dimensional test
rejects a function $f: \F_q^m \to \F_q$ with probability at least $ \frac 1 2 \cdot \min \big \{ q^{-2\ell},  \dist(f, C^{\ell \nearrow m})\big\}$. In particular, $C^{\ell \nearrow m}$ is $\left(\ell,\frac  {q^{-2\ell}} {2}\right)$-testable.
\end{theorem}

\begin{definition}[Robust testing of lifted codes]
Let $C \subseteq \Func{\ell}$ be an affine-invariant linear  code, and $m\geq k \geq \ell$.  The $m$-dimensional lift $C^{\ell \nearrow m}$ is \emph{ $(k,\alpha)$-robust} if for every $f: \F_q^m \to \F_q$ it holds that
\[\E_K\big[\dist(f|_K, C^{\ell \nearrow k})\big] \geq \alpha \cdot \dist(f, C^{\ell \nearrow m})\,,\]
where the expectation is over a uniform random $k$-dimensional affine subspace $K \subseteq \F_q^m$.
\end{definition}

We note the following easy implications.
\begin{proposition}\label[proposition]{clm:easy-implic}
Let $C \subseteq \Func{\ell}$ be an affine-invariant linear  code, and $m\geq k \geq \ell$. Then the following hold.
\begin{enumerate}
\item If $C^{\ell \nearrow m}$ is $(k,\alpha)$-robust then it is $(k,\alpha)$-testable.
\item If $C^{\ell \nearrow m}$ is $(k,\alpha)$-testable then it is $(k,\alpha \cdot q^{-k})$-robust.
\item If $C^{\ell \nearrow m}$ is $(k,\alpha)$-testable then it is $(r,\alpha)$-testable for any $k \leq r \leq m$.
\item If $C^{\ell \nearrow m}$ is $(k,\alpha)$-robust then it is $(r,\alpha)$-robust for any $k \leq r \leq m$.
\end{enumerate}
\end{proposition}

\begin{proof}
Part (1) follows since $f|_K \notin C^{\ell \nearrow k}$ whenever $\dist(f|_K, C^{\ell \nearrow k})>0$, while Part (2) follows since $\dist(f|_K, C^{\ell \nearrow k}) \geq q^{-k}$ whenever $f|_K \notin C^{\ell \nearrow k}$.

Part (3) follows by observing that for a uniform random $r$-dimensional affine subspace $R$,
\begin{align*}
\Pr_R\big[f|_R \notin C^{\ell \nearrow r}\big] & =  \E_R\left[\mathbbm{1}_{f|_R \notin C^{\ell \nearrow r}}\right] \\
 & \geq  \E_R \left[  \Pr_{K\subseteq R} \big[  f|_K \notin C^{\ell \nearrow k} \big]\right] \\
 & =  \Pr_K\big[f|_K \notin C^{\ell \nearrow k}\big]\,,
\end{align*}
where the inequality follows since $f|_R \in C^{\ell \nearrow r}$ implies that $ f|_K \in C^{\ell \nearrow k}$ for any $K$.

Finally, Part (4) follows by letting $f_R$ be the codeword in $C^{\ell \nearrow r}$ that is closest to $f|_R$, and noting that
\begin{align*}
  \E_R\big[\dist(f|_R, C^{\ell \nearrow r})\big]
  & =  \E_R\big[\dist(f|_R, f_R)\big] \\
& =  \E_R \left[ \E_{ K \subseteq R}\big[\dist(f|_K, f_R|_K)\big]\right] \\
 & \geq  \E_R \left[ \E_{ K \subseteq R}\big[\dist(f|_K, C^{\ell \nearrow k})\big]\right] \\
 & =  \E_K\big[\dist(f|_K, C^{\ell \nearrow k})\big]\,,
\end{align*}
where the inequality follows since $f_R|_K \in C^{\ell \nearrow k}$ for any $K$.
\end{proof}

We now turn to the definition of agreement testing. The agreement testers we consider are testers that for $t<s$, sample a uniform random $t$-dimensional affine subspace $T$, and a pair of uniform random $s$-dimensional affine subspaces $S,S'$ containing $T$, and accept if and only if $f_S, f_{S'}$ agree on $T$.

For a code $C \subseteq \{\F_q^m \to \F_q\}$, we let $C(s)$ be the code containing all collections $(f_S)_S$ of partial assignments to $s$-dimensional affine subspaces  that are consistent with some global codeword $g \in C$, formally,
\[ C(s):=\big\{ (f_S)_S \mid \text{$\exists \; g \in C$ such that  $g|_S = f_S$ for any $s$-dimensional affine subspace $S$}  \big\}\,. \]
For a pair of collections $(f_S)_S, (g_S)_S$ of partial assignments to $s$-dimensional affine subspaces,  we denote by $\dist((f_S)_S, (g_S)_S)$ the fraction of $s$-dimensional affine subspaces $S$ for which $f_S \neq g_S$, and we define $\dist((f_S)_S, C(s))$ accordingly.

\begin{definition}[Agreement testing of lifted codes]
Let $C \subseteq \Func{\ell}$ be an affine-invariant linear  code, and $m \geq s > t \geq \ell$. The $m$-dimensional lift $C^{\ell \nearrow m}$ is \emph{$(s,t,\alpha)$-agreement testable} if for every collection $(f_S)_S$ where $f_S \in C^{\ell \nearrow s}$ for every $s$-dimensional affine subspace $S$  it holds that
\[\Pr_{T,\; S \supseteq T, \;S' \supseteq T}[f_S|_T \neq f_{S'}|_T] \geq \alpha \cdot \dist\big((f_S)_S, C^{\ell \nearrow m}(s)\big)\,,\]
where the probability is over a uniform random $t$-dimensional affine subspace $T \subseteq \F_q^m$, and uniform random $s$-dimensional affine subspaces $S,S'$ containing $T$.
\end{definition}

\subsection{Sampling properties of affine subspaces}

%

Let $\cal T$ be a collection of affine subspaces $T \subseteq \F_q^m$ of dimension $t$, and let $\cal S$ be a collection of affine subspaces
$S \subseteq \F_q^m$ of dimension $s$, %
where %
$t<s$. Let $\calI (\cal T, \cal S)$ denote the bipartite \emph{inclusion graph}
whose left side %
is $\calT$ and  %
%
right side is $\calS$  %
and
%
%
$T\in\calT$ and $S\in\calS$ are adjacent   %
if and only if $T \subseteq S$. 
For $\alpha, \beta \in (0,1)$, we say that $\calI (\cal T, \cal S)$ is 
 an \emph{$(\alpha, \beta)$-hitting sampler} if for every subset $A \subseteq \cal T$ with $|A| \geq \alpha |\calT|$, it holds that
 $|N(A)| \geq (1-\beta)  |\cal S|$, where $N(A)$ denotes the set of neighbors of $A$ in $\cal S$. 
 
 First we note that in the case where $\cal T$ is the collection of all points in $\F_q^m$ (\ie, $0$-dimensional affine subspaces), and $\cal S$ is the collection of all $s$-dimensional affine subspaces, the hitting sampler property is an immediate consequence of the fact that points in a random affine subspace are uniformly distributed and pairwise independent, and Chebyshev's inequality.

 
\begin{lemma}\label{lem:sampling-1-s}%
Let $0<s<m$, let $\cal T$ be the collection of all points in $\F_q^m$, and let $\cal S$ be the collection of all $s$-dimensional affine subspaces of $\F_q^m$. Then for any $\alpha >0$, 
 the inclusion graph $\calI(\cal T, \cal S)$ is an $(\alpha, \frac{q^{-s}} {\alpha})$-hitting sampler.
 \end{lemma}
 
The following lemma from \cite{ImpagliazzoJKW2010}
 is an extension for the case where for a fixed affine subspace $R \subseteq \F_q^m$ of dimension $r<s$, $\mathcal T$ is the collection of all points in $\F_q^m \setminus R$, and $\mathcal S$ is the collection of all $s$-dimensional affine subspaces containing $R$.
  
 \begin{lemma}[{\cite[Lemma 2.12]{ImpagliazzoJKW2010}}]\label[lemma]{lem:sampling-1-r-s}
 Let $0<r <s<m$, and let $R \subseteq \F_q^m$ be an affine subspace of dimension $r$.
  Let $\cal T$ be the collection of all points in $\F_q^m \setminus R$, and let $\cal S$ be the collection of all $s$-dimensional affine subspaces in $\F_q^m$ containing $R$. Then for any $\alpha >0$, 
 the inclusion graph $\calI(\cal T, \cal S)$ is an $(\alpha, \frac{4 q^{-(s-r-2)}} {\alpha} )$-hitting sampler.
 \end{lemma}
 
 \begin{remark}
\cite[Lemma 2.11]{ImpagliazzoJKW2010} is only stated for the case where $s = \frac m 2$, while \cite[Lemma 2.12]{ImpagliazzoJKW2010}
is only stated for the case where $r = \frac{s} {2}$. However, it can easily be seen that the proof implies the bounds we state for general values of $0<r<s<m$.
 \end{remark}

Finally, we cite the following lemma, %
due to Impagliazzo, Kabanets and Wigderson \cite{ImpagliazzoKW2012}, which deals with the case where $\cal T$ is the collection of all $t$-dimensional \emph{linear} subspaces, and $\cal S$ is the collection of  all $s$-dimensional \emph{linear} subspaces.

\begin{lemma}[{\cite[Lemma 2.3]{ImpagliazzoKW2012}}]\label[lemma]{lem:sampling-t-s}
Let $0<t<s<m$,  let $\cal T$ be the collection of all $t$-dimensional linear subspaces in $\F_q^m$, and let $\cal S$ be the collection of all $s$-dimensional linear subspaces in $\F_q^m$. Then for any $\alpha >0$, 
 the inclusion graph $\calI(\cal T, \cal S)$ is an $(\alpha, \frac{18 q^{-(s-t)}} {\alpha})$-hitting sampler.
\end{lemma}

\begin{remark}
\cite[Lemma 2.3]{ImpagliazzoKW2012} gives a 
%
%
%
worse bound for the case that $\calI(\cal T, \cal S)$ is an 
\emph{averaging sampler}.\footnote{ For $\alpha, \beta \in (0,1)$, we say that $\calI (\cal T, \cal S)$ is 
 an \emph{$(\alpha, \beta)$-averaging sampler} if for every subset $A \subseteq \cal T$ with $|A| \geq \alpha |\calT|$, 
for at least a $(1-\beta)$-fraction of the vertices $S \in \cal S$ it holds that $\left|\frac{|A \cap  N(S)|} {\mathrm{deg}(S)} - \alpha \right| \leq \frac \alpha 2$. Note that an $(\alpha, \beta)$-averaging sampler is in particular an $(\alpha, \beta)$-hitting sampler.} However, the proof first gives the bound we state, and only then shows how to deduce the worse bound for the case of an averaging sampler.
\end{remark}

We further observe that \cref{lem:sampling-t-s} gives a similar bound for \emph{affine} subspaces.

\begin{corollary}\label[corollary]{cor:sampling-t-s}
Let $0<t<s$,  let $\cal T$ be the collection of all $t$-dimensional affine subspaces in $\F_q^m$, and let $\cal S$ be the collection of all $s$-dimensional affine subspaces in $\F_q^m$. Then for any $\alpha >0$, 
 the inclusion graph $\calI(\cal T, \cal S)$ is an $(\alpha, \frac{18 q^{-(s-t-1)}} {\alpha})$-hitting sampler.
\end{corollary}

\begin{proof}
  Let $A \subseteq \cal T$ be a collection of $t$-dimensional affine subspaces in $\F_q^m$ so that $|A| \geq \alpha |\cal T|$. Let $S \in \calS$ be a uniformly random $s$-dimensional affine subspace. Then $S$ can be sampled by first picking a random affine shift $v \in \F_q^m$, then picking a random $s$-dimensional linear subspace $V \subseteq \F_q^m$ so that $v \notin V \setminus \{0\}$, and finally letting $S = v+V$. As for a fixed vector $v \in \F_q^m$, and a random $s$-dimensional linear subspace $V \subseteq \F_q^m$, we have that $v \in V$ with probability at most $q^{-(m-s)}$, which is negligible in our setting, we may assume that $V$ is a completely random $s$-dimensional linear subspace. Next observe that $v +V$ contains a $t$-dimensional affine subspace $T = w+W \in A$ if $V$ contains $\mathrm{Span}\{w-v,W\}$. The proof is completed by noting that, conditioned on choosing $v$, by \cref{lem:sampling-t-s}, $V$ contains $\mathrm{Span}\{w-v,W\}$ for some $T = w+W \in A$ with probability at least
%
$1 - \left(18 \cdot q^{-(s-t-1)}\right)/\alpha$.  %
\end{proof}

%

\section{From agreement testing to robust testing}\label{sec:agreement-to-robust}

In this section we prove the following lemma, %
showing the agreement testing to robust testing implication.

\begin{lemma}[Agreement testing implies robust testing]\label[lemma]{agreement-to-robust}
Let $C \subseteq\{\F_{q}^{\ell}\to\F_{q}\}$ be an affine-invariant linear  code, and $m \geq s >t \geq \ell$. Suppose that $C^{\ell \nearrow m}$ is $(s,t,\alpha)$-agreement testable, and let $\delta:=\mindist(C^{\ell \nearrow t})$. Then $C^{\ell \nearrow m}$ is $\big(s, \Omega(\alpha\delta) \big)$-robust.
\end{lemma}

\begin{proof}
For simplicity of notation, in what follows we let $T, S$ denote the random variables obtained by sampling a uniform random affine subspace of dimension $t, s$ respectively.
Suppose that $f:\F_q^m \to \F_q$ has $\E_{S}[\dist(f|_S, C^{\ell \nearrow s})] \leq \epsilon$.
Our goal is to find a codeword $g \in C^{\ell \nearrow m}$ such that
$\dist(f,g)\leq O(\epsilon / (\alpha \delta))$.

The proof proceeds as follows.
We would like to apply our assumption on agreement testability, and towards this, we create an instance $(f_S)_S$ for the agreement tester by letting $f_S$ be the codeword in $C^{\ell \nearrow s}$ that is closest to $f|_S$. We then use the fact that $f_S$ is typically close to $f|_S$, together with the
fact that $t$-dimensional affine subspaces sample well inside $s$-dimensional affine subspaces, and the assumption that $C$
has large distance  %
on $t$-dimensional affine subspaces
of the domain $\F_q^s$,   %
%
to show that  
$\Pr_{T,\;S\supseteq T, \; S' \supseteq T}[f_S|_T \neq f_{S'}|_T]$ is small.
 Agreement testability then gives a codeword $g \in C^{\ell \nearrow m}$ that is consistent with most $f_S$, and using the fact that $s$-dimensional affine subspaces sample well inside $\F_q^m$ this implies in turn that $\dist(f,g)$ is small. Details follow.


We begin by showing that
$\Pr_{T, \; S \supseteq T, \; S' \supseteq T}[f_S|_T \neq f_{ S'}|_T]$  is small.
Recall first that $\E_S[\dist(f|_S, f_S)] =  \E_{S}[\dist(f|_S, C^{\ell \nearrow s})]\leq  \epsilon$. Next observe that for a fixed $s$-dimensional affine subspace $S$, any point in a uniform random $t$-dimensional affine subspace contained in $S$ is uniform in $S$. Thus we also have that
$\E_{S,\; T \subseteq S}[\dist(f|_T, f_S|_T)] \leq \epsilon,$ and consequently
\begin{align*}
 \Pr_{T, \; S \supseteq T, \; S' \supseteq T}\big[\dist(f_S|_T, f_{ S'}|_T) \geq \delta\big]
& \leq  2 \cdot \Pr_{T,\; S \supseteq T}\big[\dist(f|_T, f_S|_T) \geq \delta/2\big] \\
& =  2 \cdot \Pr_{S,\; T \subseteq S}\big[\dist(f|_T, f_S|_T) \geq \delta/2\big] \\
& \leq  \frac{4\epsilon} {\delta}\,.
\end{align*}
But since $f_S|_T$, $f_{S'}|_T$ are both codewords of $C^{\ell \nearrow t}$, a code of %
minimum   %
distance $\delta$, the above implies in turn that $\Pr_{T, \; S \supseteq T, \; S' \supseteq T}[f_S|_T \neq f_{ S'}|_T] \leq \frac{4\epsilon} {\delta}.$

Our assumption on agreement testability now gives a codeword $g \in C^{\ell \nearrow m}$ that has $\Pr_S[g|_S \neq f_S] \leq 4\epsilon/(\alpha \delta).$ But since any point in a uniform random $s$-dimensional affine subspace is uniform in $\F_q^m$ this gives in turn that
\[\dist(f,g) =\E_S\big[\dist(f|_S, g|_S)\big] \leq
  \E_S\big[\dist(f|_S, f_S)\big] + \E_S\big[\dist(f_S, g|_S)\big] \leq
  \epsilon +\frac{4\epsilon} {\alpha \delta} \leq \frac {5 \epsilon}
  {\alpha \delta} \,. \qedhere \]

\end{proof}



\section{From local testing to agreement testing}\label{sec:testing-to-agreement}

In this section we prove the following lemma that gives the local testing to agreement testing implication.

\begin{lemma}[Local testing implies agreement testing]\label[lemma]{testing-to-agreement} Let $C \subseteq\{\F_{q}^{\ell}\to\F_{q}\}$ be an affine-invariant linear  code,
and $m \geq k \geq \ell$. Suppose that
$C^{\ell \nearrow m}$ is $(k,\alpha)$-testable, and  let  $\delta:= \min_{ k \leq r \leq m}\mindist(C^{\ell \nearrow r})$.
 Then $C^{\ell \nearrow m}$ is $(2k+2+\log_q(4/\delta), k+1, \Omega(\alpha \cdot \delta^2))$-agreement testable.
\end{lemma}

\paragraph{Proof outline.}
For simplicity of notation, in what follows let $s:=2k+2+\log_q(4/\delta)$ and $t:=k+1$.
We let both $S$, $S'$ ($T$, $T'$ and $K$, $K'$, resp.) denote random variables obtained by sampling a uniform random affine subspace of dimension $s$ ($t$, $k$, resp.).

Let  $(f_S)_S$ be a collection of partial assignments such that  $f_S \in C^{\ell \nearrow s}$ for every $S$, and
\begin{equation}\label{eq:assumption}
\Pr_{T, \; S \supseteq T,\;S' \supseteq T}[f_S|_T \neq f_{S'}|_T] \leq \epsilon\,.
\end{equation}
Our goal is to find a global codeword $g \in C^{\ell \nearrow m}$ that has
\begin{equation}\label{eq:conclusion}
 \Pr_S[g|_S \neq f_S] \leq O\left(\frac{\epsilon} {\alpha \cdot \delta^2}\right)\,.
 \end{equation}

We find the codeword $g$ in three stages.
\begin{enumerate}
  \item In the initial stage (\cref{subsec:init}) we define, for any $k$-dimensional affine subspace $K$, a `most popular function' %
$\pl_K: \F_q^k \to \F_q$ by choosing the most common value among $f_S|_K$ going over all $S\supseteq K$. We show that for a typical $K$, this function is obtained with an overwhelming plurality of $1-O(\epsilon)$.
    %
  \item In the ``local structure'' stage (\cref{subsec:local-struct}) we define, for any $k$-dimensional affine subspace $K$, a function $g_K: \F_q^m \to \F_q$ by letting $g_K(x)$ be the most common `vote' %
among all $f_S$ that contain both $K$ and $x$ and agree with $\pl_K$ on $K$. We then show that for a typical $K$, $g_K$ is close to some codeword $h_K \in C^{\ell \nearrow m}$, and moreover $h_K|_S = f_S$ for most $S$ containing $K$.
  %
\item In the ``global structure" stage (\cref{subsec:global-struct}) we show that there exists $\hat K$ for which $h_{\hat K}|_S = f_S$  for most $S$ (not necessarily containing $\hat K$). We can then set our
`global function' %
$g$ to be equal to $h_{\hat K}$.
\end{enumerate}

\subsection{Initial stage}\label{subsec:init}

For any $k$-dimensional affine subspace $K$, we let $\pl_K: \F_q^k \to \F_q$ denote the most common value among $f_S|_K$ for $S$ containing $K$. That is,
\[\pl_K := \plurality_{S\supseteq K}\{f_S|_K \}\,.\]
Next we use our assumption \eqref{eq:assumption} to show that for a typical $K$, the function $\pl_K$ is obtained with overwhelming plurality.

\begin{lemma}\label[lemma]{lem:init}
  \[\E_K\bigg[\Pr_{S\supseteq K} \big[ f_S|_K \neq \pl_K \big] \bigg]\le 2\eps\,.\]
\end{lemma}
\begin{proof}
Since the collision probability lower bounds the probability of hitting the most common value, it suffices to show that
\begin{equation}\label{eq:assumption-k}
\Pr_{K, \; S \supseteq K,\;S' \supseteq K}\big[f_S|_K \neq f_{S'}|_K\big] \leq 2\epsilon\,.
\end{equation}

Clearly if $t=k$ we would be done by \eqref{eq:assumption}, so the whole point is to show the same for $t>k$. We describe a distribution on triples $(S_1,S',S_2)$ such that $(S_1,S_2)$ are distributed as in \eqref{eq:assumption-k} but the pairs $(S_1,S')$ and $(S',S_2)$ are distributed as in \eqref{eq:assumption}:
\begin{enumerate}
  \item Choose a uniform random $k$-dimensional affine subspace $K$.
  \item Choose a pair of uniform random $t$-dimensional affine subspaces $T_1,T_2$ containing $K$.
    \item For $i=1,2$, choose a uniform random $s$-dimensional affine subspace $S_i$ containing $T_i$.
  \item Choose a uniform random $s$-dimensional affine subspace $S'$ containing $T_1 \cup T_2$ (this can be done since $t=k+1$ and $s\geq k+2$).
\end{enumerate}

One can check that indeed $K,S_1,S_2$ are distributed as in \eqref{eq:assumption-k}  while  $T_i,S_i,S'$ are distributed as in \eqref{eq:assumption}. Thus by our assumption \eqref{eq:assumption},
\[ \Pr_{K, \; S_1 \supseteq K,\;S_2 \supseteq K}\big[f_{S_1}|_K \neq
  f_{S_2}|_K\big] \le   \Pr_{T_1, \; S_1 \supseteq T_1,\;S' \supseteq
    T_1}\big[ f_{S_1}|_{T_1}\neq f_{S'}|_{T_1}\big] + \Pr_{T_2, \; S'
    \supseteq T_2,\;S_2 \supseteq T_2}\big[f_{S'}|_{T_2}\neq
  f_{S_2}|_{T_2}\big] \le 2\eps\,. \qedhere \]
\end{proof}
%
%

\subsection{Local structure}\label{subsec:local-struct}

Next we define for every $k$-dimensional affine subspace $K$ the function $g_K: \F_q^m \to \F_q$. As described above,  for every $x \in \F_q^m$, we let $g_K(x)$ be the most common value among $f_S(x)$ for $S$ that contain both $K$ and $x$ and agree with $\pl_K$ on $K$. That is,
\[g_K(x) : = \plurality_{S \supseteq K \cup \{x\},\; f_S|_K = \pl_K} \{f_S(x)\}\,.\]

Next we would like to show that for a typical $K$, $g_K$ is close to some codeword $h_K \in C^{\ell \nearrow m}$, and additionally $h_K|_S = f_S$ for most $S$ containing $K$. We show these in three steps:
\begin{enumerate}
  \item Boosting step (\cref{lem:local-overwhelm-plur}): In this step we
show that for typical $K, x$, the plurality in the definition of $g_K(x)$ is obtained with overwhelming probability.
  \item LTC step (\cref{lem:local-close-to-codeword}): In this step we use the previous step to show that for a typical $g_K$,  for most $K'$, $g_K|_{K'}$ agrees with some $f_S$ on $K'$, and therefore is  a codeword of $C^{\ell \nearrow k}$. By local testability assumption this implies in turn that such $g_K$ is close to being in the code $C^{\ell \nearrow m}$, and we denote by $h_K\in C^{\ell \nearrow m}$ the closest codeword to $g_K$.
  \item Agreement step (\cref{lem:local-close-to-agree}): In this step we
  show that a typical $h_K$ agrees with most $f_S$ for $S \supseteq K$.
\end{enumerate}


We start with the boosting step, showing that for typical $K, x$, the plurality in the definition of $g_K(x)$ is obtained with overwhelming probability. Intuitively, this follows  by the assumption that the code 
has large distance %
on $t$-dimensional affine subspaces
of the domain $\F_q^s$,   %
%
together with the fact that $k$-dimensional affine subspaces sample well inside $t$-dimensional affine subspaces, which imply that if a pair of $f_S$ agree on $K$ then they must typically also agree on their whole intersection.

\begin{lemma}[Boosting step]\label[lemma]{lem:local-overwhelm-plur}
\[\mathbb{E}_{K,\; x \notin K} \left[\Pr_{S \supseteq K \cup \{x\}} \left[f_S(x) \neq g_K(x) \bigm| f_S|_K = \pl_K \right] \right]\leq  O\left(  q^{-k} \cdot \frac{\epsilon} {  \delta \cdot (1-4\epsilon) }\right)\,.\]
\end{lemma}

\begin{proof}
Since the collision probability lower bounds the probability of hitting the most common value, it suffices to show that
\[ \Pr_{K,\; x \notin K, \;S \supseteq K \cup \{x\}, \;S' \supseteq K \cup \{x\}} \big[f_S(x) \neq f_{S'}(x) \mid f_S|_K =  f_{S'}|_K = \pl_K \big]  \leq O\left(  q^{-k} \cdot \frac{\epsilon} {\delta \cdot (1-4\epsilon)}\right)\,.\]
Now we have that
\begin{align*}
 &\Pr_{K, \;x\notin K,\;S \supseteq K \cup \{x\}, \;S' \supseteq K \cup \{x\}} \big[f_S(x) \neq f_{S'}(x) \mid f_S|_K =  f_{S'}|_K = \pl_K \big] \\
&\leq  \Pr_{K, \;T\supseteq K,\;S \supseteq T, \;S' \supseteq T} \big[f_S|_{T} \neq f_{S'}|_{T} \mid f_S|_K =  f_{S'}|_K = \pl_K \big] \\
& = \frac{ \Pr_{T, \;S \supseteq T, \;S' \supseteq T, \;K\subseteq T} \big[ f_S|_K =  f_{S'}|_K = \pl_{K}  \mid f_S|_{T} \neq f_{S'}|_{T} \big]  \cdot \Pr_{T,\;S \supseteq T, \;S' \supseteq T} \big[  f_S|_{T} \neq f_{S'}|_{T} \big]} {\Pr_{K, \;T\supseteq K,\;S \supseteq T, \;S' \supseteq T} \big[ f_S|_K =  f_{S'}|_{K}  = \pl_K \big] }\,, \\
\end{align*}
where the inequality follows recalling that $t=k+1$, and so $K$ and $x$ span a random $t$-dimensional subspace.
Next we bound each of the terms above.

By our assumption \eqref{eq:assumption},  the right hand term in the numerator is %
not greater than   %
$\epsilon$. To bound the denominator note that by \cref{lem:init},
\begin{equation}\label{eq:rhs-d}
\Pr_{K, \;T\supseteq K,\;S \supseteq T, \;S' \supseteq T} \left[ f_S|_K =  f_{S'}|_{K}  = \pl_K \right]  \geq  1- 2\cdot \Pr_{K, \;S \supseteq K} \left[ f_S|_K \neq \pl_K \right] \geq 1-4\epsilon\,.
\end{equation}

%
To bound the left hand term in the numerator note first that since $f_{S}, f_{S'} $ are both codewords of  $C^{\ell \nearrow s}$ then $f_{S}|_T, f_{S'}|_T$ are distinct codewords in $C^{\ell \nearrow t}$, and so
 $\dist(f_S|_T, f_{S'}|_T) \geq \delta$.
We now apply \cref{lem:sampling-1-s} on the graph $\calI(T, \cal K)$, where the ambient space is $T$,  $T$ contains all points in $T$, and $\cal K$ contains all $k$-dimensional subspaces in $T$. By \cref{lem:sampling-1-s}, the graph $\calI ( T, \cal K)$ is a $(\delta, \frac {q^{-k}} {\delta})$-hitting sampler, and so taking
$ A'=\sett{x\in T}{f_S(x)\neq f_{S'}(x)}$
we deduce that at most a $\frac{q^{-k}} \delta$-fraction of $K$ can miss $A$
altogether. Thus,
 \begin{equation}\label{eq:lhs}
 \Pr_{T, \;S \supseteq T, \;S' \supseteq T, \;K \subseteq T} \left[ f_S|_K =  f_{S'}|_K  \bigm| f_S|_{T} \neq f_{S'}|_{T} \right]  \leq \frac{q^{-k}} {\delta} \,.
 \end{equation}

 The final bound is obtained by combining the bounds in \eqref{eq:assumption}, \eqref{eq:rhs-d}, and \eqref{eq:lhs}.
\end{proof}


Next we use the assumption on local testability to
show that for a typical $K$, $g_K$ is close to being a codeword of $C^{\ell \nearrow m}$.

\begin{lemma}[LTC step]\label[lemma]{lem:local-close-to-codeword}
\[\mathbb{E}_{K}\left[\dist\big(g_K,C^{\ell \nearrow m}\big)\right] \leq  O\left(  \frac{\epsilon} {\alpha \cdot \delta}\right)\,.\]
\end{lemma}

\begin{proof}
To apply our assumption on local testability we first show that $g_K|_{K'}$ is typically a codeword of $C^{\ell \nearrow k}$. For this, first observe that if $g_K|_{K'}$ is not a codeword of $C^{\ell \nearrow k}$ then $g_K|_{K'} \neq f_S|_{K'}$ for all $S$ (since $f_S \in C^{\ell \nearrow s}$ and so $f_S|_{K'} \in C^{\ell \nearrow k}$). Thus we have
\begin{align*}
 \E_{K,\;K'}\left[\mathbbm{1}_{g_K|_{K'} \notin C^{\ell \nearrow k}}\right]
&  \leq   \E_{K,{\;K'}}\left[\Pr_{ S \supseteq K \cup K'} \left[g_K|_{K'} \neq f_S|_{K'}  \right] \right] \\
 & \leq  \E_{K,{\;K'}}\left[\Pr_{ S \supseteq K \cup K'} \left[g_K|_{K'} \neq f_S|_{K'} \bigm|  f_S|_K = \pl_K \right] \right] + \Pr_{K,\;S\supseteq K}[ f_S|_K \neq \pl_K]\,.
 \end{align*}

 We claim that the above expression is at most $O(\epsilon/\delta)$.
 To see this note first that by \cref{lem:init} the right hand term is at most $2\epsilon$ . To bound the left hand term, note that since each individual point in $K'$ is uniformly distributed in $\F_q^m$ this term is
%
not greater than  %
\[q^k \cdot \E_{K,\;x}\left[\Pr_{ S \supseteq K \cup \{x\}} \left[g_K(x) \neq f_S(x) \bigm|  f_S|_K = \pl_K \right] \right]\,,\] which is in turn at most $O(\epsilon/\delta)$  by \cref{lem:local-overwhelm-plur} (noting that the probability in the above expression is zero whenever $x \in K$).

Finally, for any $k$-dimensional affine subspace $K$, let $\epsilon_{K}:= \Pr_{K'}\left[g_K|_{K'} \notin C^{\ell \nearrow k}\right]$. Then on the one hand
 $\E_K[\epsilon_K]  = \E_{K,\;K'}\left[\mathbbm{1}_{g_K|_{K'} \notin C^{\ell \nearrow k}}\right]  \leq O(\epsilon/\delta)$, and on the other hand $\dist(g_K, C^{\ell \nearrow m}) \leq \epsilon_K/\alpha$ for any $K$
  by assumption that $C^{\ell \nearrow m}$ is  $(k,\alpha)$-testable. We  conclude that
\[ \mathbb{E}_{K}\left[\dist\big(g_K,C^{\ell \nearrow m}\big)\right]
\leq \E_{K}\left[\frac{\epsilon_K} {\alpha}\right ] \leq O\left(
  \frac{\epsilon} {\alpha \cdot \delta}\right)\,. \qedhere \]
\end{proof}

For any $k$-dimensional affine subspace $K$, let $h_K \in C^{\ell \nearrow m}$ be the codeword that is closest to $g_K$. Then by the above lemma,
\[\mathbb{E}_{K}\big[\dist(g_K,h_K)\big] \leq   O\left(\frac{\epsilon} {\alpha \cdot \delta}\right)\,.\]
The following lemma says that for a typical $K$, we have that $h_K|_S = f_S$ for most $S$ containing $K$, which follows by the fact that $s$-dimensional affine subspaces sample well inside $\F_q^m$ and by assumption that
the code
has large distance %
on $s$-dimensional affine subspaces
of the domain $\F_q^m$.   %
%

\begin{lemma}[Agreement step]\label[lemma]{lem:local-close-to-agree}
\[\E_{K}\left[\Pr_{\;S \supseteq K}\big[ h_K|_S \neq f_S\big]\right] \leq O\left(\frac{\epsilon} {\alpha \cdot \delta^2}\right)\,.\]
\end{lemma}

\begin{proof}
  We show that for typical $S \supseteq K$, on the one hand, by \cref{lem:local-close-to-codeword} and the fact that $s$-dimensional affine subspaces sample well inside $\F_q^m$, $\dist(h_K|_S, g_K|_S)$ is small, and on the other hand, by \cref{lem:local-overwhelm-plur}, $\dist(g_K|_S, f_S)$ is small. We then conclude by triangle inequality that $\dist(h_K|_S, f_S)$ is  small, which implies in turn that $h_K|_S = f_S$ by assumption that the code
has large distance  %
on $s$-dimensional subspaces
of the domain $\F_q^m$.   %
%

We start by showing that $\dist(h_K|_S, g_K|_S)$ is typically small.
For this note that for a fixed $k$-dimensional affine subspace $K$, and uniform random $S$ containing $K$, each individual point in $S \setminus K$ is uniformly distributed in $\F_q^m\setminus K$. Thus we have
\[
 \E_{K,\;S \supseteq K}\big[\dist(h_{K}|_{S\setminus K},g_{K}|_{S \setminus K})\big] \leq \frac{\mathbb{E}_{K}\big[\dist(h_K,g_K)\big] } {1- q^{-(m-k)}}\leq O\left(\frac{\epsilon} {\alpha \cdot \delta}\right)\,,
\]
 where the last inequality follows by \cref{lem:local-close-to-codeword}. Markov's inequality then implies that 
 $\dist(h_{K}|_{S\setminus K},g_{K}|_{S \setminus K})$
%
$ \leq \frac{\delta} {4}$ with probability at least $1- O\left(\frac{\epsilon} {\alpha \cdot \delta^2}\right)$ over the choice of  $K$ and $S \supseteq K$. We conclude that with probability at least $1- O\left(\frac{\epsilon} {\alpha \cdot \delta^2}\right)$ over the choice of  $K$ and $S \supseteq K$, it holds that
  \begin{equation}\label{eq:g-h-close-on-T}
 \dist(h_K|_S, g_K|_S) \leq \frac{\delta} {4} + q^{-(s-k)}\leq \frac{\delta} {2}\, ,
  \end{equation}
 where the last inequality follows by choice of $s \geq k + \log_q(4/\delta)$.

Next we show that $\dist(g_K|_S, f_S)$ is  typically small.
For this note that
\begin{align}\label{eq:h-close-to-f-on-T}
\nonumber &  \E_{K,\;S \supseteq K}\big[\dist(g_{K}|_S,f_{S})\big] \\
 \nonumber  & \leq \Pr_{K,\;S \supseteq K}[ g_K|_K \neq f_S|_K]+   \E_{K,\;S \supseteq K}\big[\dist(g_{K}|_S,f_{S}) \mid g_K|_K = f_S|_K\big] \\
  \nonumber  & \leq \Pr_{K,\;S \supseteq K}[ g_K|_K \neq f_S|_K]+   \E_{K,\;S \supseteq K}\big[\dist(g_{K}|_{S\setminus K},f_{S}|_{S\setminus K}) \mid g_K|_K = f_S|_K\big] \\
  \nonumber & =  \Pr_{K,\;S \supseteq K} \left[f_S|_K \neq \pl_K \right]+\mathbb{E}_{K,\; x \notin K} \left[\Pr_{S \supseteq K \cup \{x\}} \big[g_K(x) \neq f_S(x) \bigm| f_S|_K = \pl_K \big] \right]  \\
& \leq   O\left( \frac{\epsilon} { \delta}\right)\,,
\end{align}
where the last inequality follows by  \cref{lem:init,lem:local-overwhelm-plur}. Markov's inequality then implies that 
\begin{equation}\label{eq:g-close-to-f-on-T}
 \dist(g_{K}|_S,f_S) \leq \frac{\delta} {4} 
 \end{equation}
 with probability at least $1- O\left(\frac{\epsilon} {\delta^2}\right)$ over the choice of  $K$ and $S \supseteq K$. 
 

Combining \eqref{eq:g-h-close-on-T} and \eqref{eq:h-close-to-f-on-T}, by a union bound and triangle inequality, we have that
 $\dist(h_{K}|_S,f_S) < \delta$ with probability at least $1- O\left(\frac{\epsilon} {\alpha \cdot \delta^2}\right)$ over the choice of  $K$ and $S \supseteq K$. 
Finally, since both $h_K|_{S}$ and $f_{S}$ are codewords of $C^{\ell \nearrow s}$ and $\mindist(C^{\ell \nearrow s}) \geq \delta$ we conclude that  $h_K|_{S} \neq f_{S}$ with probability at most $O\left(\frac{\epsilon} {\alpha \cdot \delta^2}\right)$  over the choice of $K$ and $S \supseteq K$.
\end{proof}


\subsection{Global structure}\label{subsec:global-struct}

We now complete the proof of \cref{testing-to-agreement} by showing that there exists a codeword $g \in C^{\ell \nearrow m}$ that agrees with most $f_S$. We start by showing that most functions $h_K$ are in fact identical, which follows by
\cref{lem:local-close-to-agree} and the fact that $s$-dimensional affine subspaces sample well inside $\F_q^m$.

\begin{lemma}\label[lemma]{lem:g-exists}
There exists a $k$-dimensional affine subspace $\hat K$ such that
\[\Pr_{K}\big[h_{K} \neq h_{\hat K}\big] \leq O\left(\frac{\epsilon} {\alpha \cdot \delta^2}\right) \,.\]
\end{lemma}

\begin{proof}
By \cref{lem:local-close-to-agree},
\[
\Pr_{K, \; K',\; S \supseteq K \cup K'}\big[h_K|_S \neq h_{K'}|_S\big] \leq 2\cdot \Pr_{K, \; S \supseteq K}\big[h_K|_S \neq f_S\big] \leq O\left(\frac{\epsilon} {\alpha \cdot \delta^2}\right)\,,
\]
and so by averaging there exists $\hat K$ such that
\[
\Pr_{K, \; S \supseteq  K \cup \hat K}\big[h_{ K}|_S \neq h_{\hat K}|_S\big] \leq  O\left(\frac{\epsilon} {\alpha \cdot \delta^2}\right)\,.
\]
Markov's inequality then implies that
\[\Pr_{S \supseteq K \cup \hat K}\big[h_{ K}|_S \neq h_{ \hat K}|_S\big]  \geq \frac{1} {2}\]
with probability at most $O\left(\frac{\epsilon} {\alpha \cdot \delta^2}\right)$ over the choice of $K$.

%
Next we claim that when $h_{K} \neq h_{\hat K} $ then $\Pr_{S \supseteq K \cup \hat K}\big[h_{ K}|_S \neq h_{ \hat K}|_S\big]  \geq \frac{1} {2},$ and so $h_{K} \neq h_{\hat K}$ with probability at most $O\left(\frac{\epsilon} {\alpha \cdot \delta^2}\right)$ over the choice of $K$. To this end, observe that if $h_{K} \neq h_{\hat K} $ then since $h_{ K}$, $h_{ \hat K}$ are both codewords of $C^{\ell \nearrow m}$ and $\mindist(C^{\ell \nearrow m}) \geq \delta$, it must hold that $\dist(h_{ K}, h_{\hat K}) \geq \delta$. If $h_K$ and $h_{\hat K}$ differ on some point of $K \cup \hat K$, then we clearly have that $\Pr_{S \supseteq K \cup \hat K}\big[h_{ K}|_S \neq h_{ \hat K}|_S\big]  =1,$ and so we are done. Hence, we may assume that $h_K$ agrees with $h_{\hat K}$ on $K \cup \hat K$, and so $\dist \left(h_{ K}|_{ \F_q^m \setminus (K \cup \hat K)}, h_{\hat K}| _{  \F_q^m \setminus (K \cup \hat K)  } \right)\geq \delta$.

We now apply \cref{lem:sampling-1-r-s} on the graph $\calI ( \F_q^m \setminus (K \cup \hat K),\cal S )$ that connects the points of $\F_q^m\setminus (K \cup \hat K)$ on the left to the $s$-dimensional affine subspaces containing $K \cup \hat K$ on the right. By \cref{lem:sampling-1-r-s}, the graph $\calI ( \F_q^m \setminus (K \cup \hat K),\cal S )$ is a $(\delta, \frac{4q^{-(s-2k-2)} } {\delta})$-hitting sampler, and so taking
%
\[
A'=\sett{x\in \F_q^m \setminus (K \cup \hat K)}{h_{ K}(x)\neq h_{\hat K}(x)},
\]
we deduce that
\[
\Pr_{S \supseteq \hat K \cup K}\big[h_{\hat K}|_S \neq h_{K}|_S\big] \geq 1- \frac{ q^{-(s-2k-2)}} {\delta} \geq \frac12\,,
\]
where the last inequality follows by assumption that $s \geq 2k +2+\log_q(2/\delta)$.

%
%
%
%
%

It now follows that $h_{K} \neq h_{\hat K}$ with probability at most $O\left(\frac{\epsilon} {\alpha \cdot \delta^2}\right)$ over the choice of $K$.
\end{proof}

We can now complete the proof of \cref{testing-to-agreement}.

\begin{proof}[Proof of \cref{testing-to-agreement}]
Set $g \in C^{\ell \nearrow m}$  to be the function $h_{\hat K}$ guaranteed  by \cref{lem:g-exists}.
By \cref{lem:local-close-to-agree,lem:g-exists},
\[
\Pr_S\big[g|_S \neq f_S\big] =\Pr_{K, \; S \supseteq K}\big[g|_S \neq f_S\big]   \leq \Pr_{K}\big[g \neq h_{K}\big] + \Pr_{K, \; S \supseteq K}\big[h_{K}|_S \neq f_S\big] \leq O\left(\frac{\epsilon} {\alpha \cdot \delta^2}\right)\,.
\]
So $g$ satisfies \eqref{eq:conclusion} as required.
\end{proof}

\section{From local testing to robust testing}\label{sec:testing-to-robust}

\subsection{Proof of \cref{thm:main}}

We can now combine  \cref{agreement-to-robust,testing-to-agreement} to prove our \cref{thm:main}, restated below, showing a transformation from local testing to robust testing.

\maint*

%
%
%
%
%
%

\begin{proof}[Proof of \cref{thm:main}]
%
%

By \cref{testing-to-agreement} we have that $C^{\ell \nearrow m}$ is
$(2k+2+\log_q(4/\delta), k+1, \Omega(\alpha \cdot \delta^2))$-agreement testable, and by \cref{agreement-to-robust} this implies in turn that $C^{\ell \nearrow m}$ is $(2k+2+\log_q(4/\delta),  \Omega(\alpha \cdot \delta^3))$-robust.
\end{proof}


\subsection{Proof of \cref{cor:main}}

We now instantiate our \cref{thm:main} with \cref{lifted-testing} to show that lifted codes are robustly testable. For this, we first observe that
one can amplify the soundness of the tester given by \cref{lifted-testing}
to a constant (independent of $q$ and $\ell$) at the cost of increasing the testing dimension to $\approx 3\ell$.

\begin{proposition}\label[proposition]{lifted-testing-ampl}
Let $C \subseteq\{\F_{q}^{\ell}\to\F_{q}\}$ be an affine-invariant linear  code, and $m \geq 3\ell+1+\log_q72$.  Then $C^{\ell \nearrow m}$ is $\left(3\ell+1+\log_q72, \Omega(1) \right)$-testable.
\end{proposition}

\begin{proof}
If the $\ell$-dimensional test rejects with probability at least $\frac 1 2 \cdot \dist(f, C^{\ell \nearrow m})$ then by Part (3) of \cref{clm:easy-implic}, the ($3\ell+\log_q4$)-dimensional test also rejects with the same probability and we are done. Otherwise, by \cref{lifted-testing}, the $\ell$-dimensional test rejects with probability at least $\frac 1 2 \cdot q^{-2\ell}$.

%
Consider the  graph $\calI( \cal S, \cal T)$ with left hand side being all $\ell$-dimensional affine subspaces of $\F_q^m$ and right hand side being all ($3\ell+1+\log_q 72$)-dimensional affine subspaces of $\F_q^m$. Next we apply \cref{cor:sampling-t-s} on the graph $\calI (\cal S, \cal T)$ with
$A'$ being the collection of all $\ell$-dimensional affine subspaces on which the $\ell$-dimensional test rejects.
Noting that the ($3\ell+1+\log_q72$)-dimensional test will reject on any neighbor of $A$ we conclude that the ($3\ell+1+\log_q72$)-dimensional test rejects with probability at least
\[1 - \frac{18q^{-(2\ell+\log_q72)}} {q^{-2\ell} /2 }=1 -
\frac{q^{-(2\ell+\log_q4)}} {q^{-2\ell} /2 } = \frac{1} {2}\,. \qedhere\]
\end{proof}

We now turn to the proof of \cref{cor:main}, restated below.

\maincor*

%
%
%
%
%
%
%

\begin{proof}
Suppose first that $\delta < 2q^{-\ell}$.
In this case by \cref{lifted-testing}, $C^{\ell \nearrow m}$ is $(\ell,\Omega(q^{-2\ell}))$-testable, and so by Part (2) of
\cref{clm:easy-implic},
 $C^{\ell \nearrow m}$ is also robustly testable using the $\ell$-dimensional test with robust soundness $\Omega\big(q^{-3\ell}\big) \geq \Omega(\delta^3)$. By Part (4) of \cref{clm:easy-implic}
it follows that the $(6\ell +4+\log_q(c/\delta))$-dimensional test also has robust soundness $\Omega(\delta^3)$.

Next assume that $\delta \geq 2q^{-\ell}$. In this case \cref{lifted-dist} gives that $\mindist(C^{\ell \nearrow r})\geq \delta/2$ for any $\ell \leq r \leq m$, and so we may apply \cref{lifted-testing-ampl} and
 \cref{thm:main} and conclude that  $C^{\ell \nearrow m}$ is $(6\ell+4+\log_q(c/\delta),  \Omega( \delta^3))$-robust for an absolute constant $c >1$.
\end{proof}



%% Theory of Computing
%% File: v018a012.bbl
%% Publication date: June 6, 2022
%% Authors: Irit Dinur, Prahladh Harsha, Tali Kaufman, and Noga Ron-Zewi
\newcommand{\prelim}{Preliminary version in\ }
% Define empty bibhead if not already defined
\providecommand{\bibhead}[1]{}
% Define tocrefpdfbookmark if not already defined
\expandafter\ifx\csname pdfbookmark\endcsname\relax%
  \providecommand{\tocrefpdfbookmark}{}
\else\providecommand{\tocrefpdfbookmark}{%
   \hypertarget{tocreferences}{}%
   \pdfbookmark[1]{References}{tocreferences}}%
\fi

\tocrefpdfbookmark
\begin{thebibliography}{10}

\bibitem{Arora1994}\bibhead{Arora1994}
{\sc Sanjeev Arora}: {\em Probabilistic Checking of Proofs and the Hardness of
  Approximation Problems}.
\newblock Ph.\,D.\ thesis, UC Berkeley, 1994.
\newblock
  \href{https://eccc.weizmann.ac.il/static/books/Probabilistic_Checking_of_Proofs_and_Hardness_of_Approximation_Problems/}{ECCC}.

\bibitem{AroraLMSS1998}\bibhead{AroraLMSS1998}
{\sc Sanjeev Arora, Carsten Lund, Rajeev Motwani, Madhu Sudan, and Mario
  Szegedy}: Proof verification and the hardness of approximation problems.
\newblock {\em J. ACM}, 45(3):501--555, 1998.
\newblock \prelim \href{https://doi.org/10.1109/SFCS.1992.267823}{FOCS'92}.
\newblock [\epfmtdoi{10.1145/278298.278306}, \epfmt{eccc}{TR98-008}]

\bibitem{AroraS1998}\bibhead{AroraS1998}
{\sc Sanjeev Arora and Shmuel Safra}: Probabilistic checking of proofs: A new
  characterization of~{NP}.
\newblock {\em J. ACM}, 45(1):70--122, 1998.
\newblock \prelim \href{https://doi.org/10.1109/SFCS.1992.267824}{FOCS'92}.
\newblock [\epfmtdoi{10.1145/273865.273901}]

\bibitem{AroraS2003}\bibhead{AroraS2003}
{\sc Sanjeev Arora and Madhu Sudan}: Improved low-degree testing and its
  applications.
\newblock {\em Combinatorica}, 23(3):365--426, 2003.
\newblock \prelim \href{https://doi.org/10.1145/258533.258642}{STOC'97}.
\newblock [\epfmtdoi{10.1007/s00493-003-0025-0}, \epfmt{eccc}{TR97-003}]

\bibitem{BabaiFLS1991}\bibhead{BabaiFLS1991}
{\sc L{\'a}szl{\'o} Babai, Lance Fortnow, Leonid~A. Levin, and Mario Szegedy}:
  Checking computations in polylogarithmic time.
\newblock In {\em Proc. 23rd STOC}, pp. 21--31. ACM Press, 1991.
\newblock [\epfmtdoi{10.1145/103418.103428}]

\bibitem{BabaiFL1991}\bibhead{BabaiFL1991}
{\sc L{\'a}szl{\'o} Babai, Lance Fortnow, and Carsten Lund}: Non-deterministic
  exponential time has two-prover interactive protocols.
\newblock {\em Comput. Complexity}, 1(1):3--40, 1991.
\newblock \prelim \href{https://doi.org/10.1109/FSCS.1990.89520}{FOCS'90}.
\newblock [\epfmtdoi{10.1007/BF01200056}]

\bibitem{BenSassonGHSV2006}\bibhead{BenSassonGHSV2006}
{\sc Eli {Ben-Sasson}, Oded Goldreich, Prahladh Harsha, Madhu Sudan, and Salil
  Vadhan}: Robust {PCP}s of proximity, shorter {PCP}s and applications to
  coding.
\newblock {\em SIAM J. Comput.}, 36(4):889--974, 2006.
\newblock \prelim \href{https://doi.org/10.1145/1007352.1007361}{STOC'04}.
\newblock [\epfmtdoi{10.1137/S0097539705446810}, \epfmt{eccc}{TR04-021}]

\bibitem{BenSassonHR2005}\bibhead{BenSassonHR2005}
{\sc Eli {Ben-Sasson}, Prahladh Harsha, and Sofya Raskhodnikova}: Some {3CNF}
  properties are hard to test.
\newblock {\em SIAM J. Comput.}, 35(1):1--21, 2005.
\newblock \prelim \href{https://doi.org/10.1145/780542.780594}{STOC'03}.
\newblock [\epfmtdoi{10.1137/S0097539704445445}]

\bibitem{BenSassonMSS2011}\bibhead{BenSassonMSS2011}
{\sc Eli Ben{-}Sasson, Ghid Maatouk, Amir Shpilka, and Madhu Sudan}: Symmetric
  {LDPC} codes are not necessarily locally testable.
\newblock In {\em Proc. 26th IEEE Conf. on Comput. Complexity (CCC'11)}, pp.
  55--65. IEEE Comp. Soc., 2011.
\newblock [\epfmtdoi{10.1109/CCC.2011.14}, \epfmt{eccc}{TR10-199}]

\bibitem{BenSassonS2006}\bibhead{BenSassonS2006}
{\sc Eli Ben{-}Sasson and Madhu Sudan}: Robust locally testable codes and
  products of codes.
\newblock {\em Random Struct. Algor.}, 28(4):387--402, 2006.
\newblock \prelim
  \href{https://doi.org/10.1007/978-3-540-27821-4\_26}{RANDOM'04}.
\newblock [\epfmtdoi{10.1002/rsa.20120}, \epfmt{arxiv}{cs/0408066},
  \epfmt{eccc}{TR04-046}]

\bibitem{BenSassonV2015}\bibhead{BenSassonV2015}
{\sc Eli Ben{-}Sasson and Michael Viderman}: Composition of semi-{LTC}s by
  two-wise tensor products.
\newblock {\em Comput. Complexity}, 24(3):601--643, 2015.
\newblock \prelim
  \href{https://doi.org/10.1007/978-3-642-03685-9\_29}{RANDOM'09}.
\newblock [\epfmtdoi{10.1007/s00037-013-0074-8}, \epfmt{eccc}{TR11-070}]

\bibitem{BlumLR1993}\bibhead{BlumLR1993}
{\sc Manuel Blum, Michael Luby, and Ronitt Rubinfeld}: Self-testing/correcting
  with applications to numerical problems.
\newblock {\em J. Comput. System Sci.}, 47(3):549--595, 1993.
\newblock \prelim \href{https://doi.org/10.1145/100216.100225}{STOC'90}.
\newblock [\epfmtdoi{10.1016/0022-0000(93)90044-W}]

\bibitem{DiksteinD2019}\bibhead{DiksteinD2019}
{\sc Yotam Dikstein and Irit Dinur}: Agreement testing theorems on layered set
  systems.
\newblock In {\em Proc. 60th FOCS}, pp. 1495--1524. IEEE Comp. Soc., 2019.
\newblock [\epfmtdoi{10.1109/FOCS.2019.00088}, \epfmt{arxiv}{1909.00638},
  \epfmt{eccc}{TR19-112}]

\bibitem{DinurG2008}\bibhead{DinurG2008}
{\sc Irit Dinur and Elazar Goldenberg}: Locally testing direct product in the
  low error range.
\newblock In {\em Proc. 49th FOCS}, pp. 613--622. IEEE Comp. Soc., 2008.
\newblock [\epfmtdoi{10.1109/FOCS.2008.26}]

\bibitem{conf-version}\bibhead{conf-version}
{\sc Irit Dinur, Prahladh Harsha, Tali Kaufman, and Noga Ron-Zewi}: From local
  to robust testing via agreement testing.
\newblock In {\em Proc. 10th Innovations in Theoret. Comp. Sci. Conf.
  (ITCS'19)}, pp. 29:1--18. Schloss Dagstuhl--Leibniz-Zentrum fuer Informatik,
  2019.
\newblock [\epfmtdoi{10.4230/LIPIcs.ITCS.2019.29}]

\bibitem{DinurK2017}\bibhead{DinurK2017}
{\sc Irit Dinur and Tali Kaufman}: High dimensional expanders imply agreement
  expanders.
\newblock In {\em Proc. 58th FOCS}, pp. 974--985. IEEE Comp. Soc., 2017.
\newblock [\epfmtdoi{10.1109/FOCS.2017.94}, \epfmt{eccc}{TR17-089}]

\bibitem{DinurN2017}\bibhead{DinurN2017}
{\sc Irit Dinur and Inbal~Livni Navon}: Exponentially small soundness for the
  direct product {Z}-test.
\newblock In {\em Proc. 32nd Comput. Complexity Conf. (CCC'17)}, pp.
  29:1--29:50. Schloss Dagstuhl--Leibniz-Zentrum fuer Informatik, 2017.
\newblock [\epfmtdoi{10.4230/LIPIcs.CCC.2017.29}, \epfmt{eccc}{TR17-096}]

\bibitem{DinurR2006}\bibhead{DinurR2006}
{\sc Irit Dinur and Omer Reingold}: Assignment testers: {T}owards a
  combinatorial proof of the {PCP} {T}heorem.
\newblock {\em SIAM J. Comput.}, 36(4):975--1024, 2006.
\newblock \prelim \href{https://doi.org/10.1109/FOCS.2004.16}{FOCS'04}.
\newblock [\epfmtdoi{10.1137/S0097539705446962}]

\bibitem{DinurS2014-dp}\bibhead{DinurS2014-dp}
{\sc Irit Dinur and David Steurer}: Direct product testing.
\newblock In {\em Proc. 29th IEEE Conf. on Comput. Complexity (CCC'14)}, pp.
  188--196. IEEE Comp. Soc., 2014.
\newblock [\epfmtdoi{10.1109/CCC.2014.27}, \epfmt{eccc}{TR13-179}]

\bibitem{FeigeGLSS1996}\bibhead{FeigeGLSS1996}
{\sc Uriel Feige, Shafi Goldwasser, L{\'a}szl{\'o} Lov{\'a}sz, Shmuel Safra,
  and Mario Szegedy}: Interactive proofs and the hardness of approximating
  cliques.
\newblock {\em J. ACM}, 43(2):268--292, March 1996.
\newblock \prelim \href{https://doi.org/10.1109/SFCS.1991.185341}{FOCS'91}.
\newblock [\epfmtdoi{10.1145/226643.226652}]

\bibitem{GemmellLRSW1991}\bibhead{GemmellLRSW1991}
{\sc Peter Gemmell, Richard~J. Lipton, Ronitt Rubinfeld, Madhu Sudan, and Avi
  Wigderson}: Self-testing/correcting for polynomials and for approximate
  functions.
\newblock In {\em Proc. 23rd STOC}, pp. 32--42. ACM Press, 1991.
\newblock [\epfmtdoi{10.1145/103418.103429}]

\bibitem{GoldreichS2000}\bibhead{GoldreichS2000}
{\sc Oded Goldreich and Shmuel Safra}: A combinatorial consistency lemma with
  application to proving the {PCP} theorem.
\newblock {\em SIAM J. Comput.}, 29(4):1132--1154, 2000.
\newblock \prelim \href{https://doi.org/10.1007/3-540-63248-4\_7}{RANDOM'97}.
\newblock [\epfmtdoi{10.1137/S0097539797315744}, \epfmt{eccc}{TR96-047}]

\bibitem{GuoHS2015}\bibhead{GuoHS2015}
{\sc Alan Guo, Elad Haramaty, and Madhu Sudan}: Robust testing of lifted codes
  with applications to low-degree testing.
\newblock In {\em Proc. 56th FOCS}, pp. 825--844. IEEE Comp. Soc., 2015.
\newblock [\epfmtdoi{10.1109/FOCS.2015.56}, \epfmt{eccc}{TR15-043}]

\bibitem{GuoKS2013}\bibhead{GuoKS2013}
{\sc Alan Guo, Swastik Kopparty, and Madhu Sudan}: New affine-invariant codes
  from lifting.
\newblock In {\em Proc. 4th Innovations in Theoret. Comp. Sci. Conf.
  (ITCS'13)}, pp. 529--540. ACM Press, 2013.
\newblock [\epfmtdoi{10.1145/2422436.2422494}, \epfmt{arxiv}{1208.5413},
  \epfmt{eccc}{TR12-149}]

\bibitem{HaramatyRS2015}\bibhead{HaramatyRS2015}
{\sc Elad Haramaty, Noga Ron{-}Zewi, and Madhu Sudan}: Absolutely sound testing
  of lifted codes.
\newblock {\em Theory of Computing}, 11(12):299--338, 2015.
\newblock \prelim
  \href{https://doi.org/10.1007/978-3-642-40328-6\_46}{RANDOM'13}.
\newblock [\epfmtdoi{10.4086/toc.2015.v011a012}, \epfmt{eccc}{TR13-030}]

\bibitem{ImpagliazzoJKW2010}\bibhead{ImpagliazzoJKW2010}
{\sc Russell Impagliazzo, Ragesh Jaiswal, Valentine Kabanets, and Avi
  Wigderson}: Uniform direct product theorems: Simplified, optimized, and
  derandomized.
\newblock {\em SIAM J. Comput.}, 39(4):1637--1665, 2010.
\newblock \prelim \href{https://doi.org/10.1145/1374376.1374460}{STOC'08}.
\newblock [\epfmtdoi{10.1137/080734030}, \epfmt{eccc}{TR08-079}]

\bibitem{ImpagliazzoKW2012}\bibhead{ImpagliazzoKW2012}
{\sc Russell Impagliazzo, Valentine Kabanets, and Avi Wigderson}: New
  direct-product testers and 2-query {PCP}s.
\newblock {\em SIAM J. Comput.}, 41(6):1722--1768, 2012.
\newblock \prelim \href{https://doi.org/10.1145/1536414.1536435}{STOC'09}.
\newblock [\epfmtdoi{10.1137/09077299X}, \epfmt{eccc}{TR09-090}]

\bibitem{KaufmanS2008}\bibhead{KaufmanS2008}
{\sc Tali Kaufman and Madhu Sudan}: Algebraic property testing: the role of
  invariance.
\newblock In {\em Proc. 40th STOC}, pp. 403--412. ACM Press, 2008.
\newblock [\epfmtdoi{10.1145/1374376.1374434}, \epfmt{eccc}{TR07-111}]

\bibitem{KoppartyMRS2017}\bibhead{KoppartyMRS2017}
{\sc Swastik Kopparty, Or~Meir, Noga Ron{-}Zewi, and Shubhangi Saraf}:
  High-rate locally correctable and locally testable codes with sub-polynomial
  query complexity.
\newblock {\em J. ACM}, 64(2):11:1--11:42, 2017.
\newblock \prelim \href{https://doi.org/10.1145/2897518.2897523}{STOC'16}.
\newblock [\epfmtdoi{10.1145/3051093}, \epfmt{arxiv}{1504.05653},
  \epfmt{eccc}{TR15-068}]

\bibitem{RazS1997}\bibhead{RazS1997}
{\sc Ran Raz and Shmuel Safra}: A sub-constant error-probability low-degree
  test, and a sub-constant error-probability {PCP} characterization of~{NP}.
\newblock In {\em Proc. 29th STOC}, pp. 475--484. ACM Press, 1997.
\newblock [\epfmtdoi{10.1145/258533.258641}]

\bibitem{RubinfeldS1996}\bibhead{RubinfeldS1996}
{\sc Ronitt Rubinfeld and Madhu Sudan}: Robust characterizations of polynomials
  with applications to program testing.
\newblock {\em SIAM J. Comput.}, 25(2):252--271, 1996.
\newblock Preliminary versions in
  \href{https://doi.org/10.1145/103418.103429}{STOC'91} and
  \href{https://dl.acm.org/doi/10.5555/139404.139410}{SODA'92}.
\newblock [\epfmtdoi{10.1137/S0097539793255151}]

\bibitem{Viderman2015}\bibhead{Viderman2015}
{\sc Michael Viderman}: A combination of testability and decodability by tensor
  products.
\newblock {\em Random Struct. Algor.}, 46(3):572--598, 2015.
\newblock \prelim
  \href{https://doi.org/10.1007/978-3-642-32512-0\_55}{RANDOM'12}.
\newblock [\epfmtdoi{10.1002/rsa.20498}, \epfmt{arxiv}{1105.5806},
  \epfmt{eccc}{TR11-087}]

\end{thebibliography}

\end{document}
